\begin{theorem}[Exact stable twist]\label{mod:lamprecht-stabletwist}
Let $F$ be a nonarchimedean local field, let $\psi$ have exact additive
conductor $n$, and let $\theta$ and $\nu$ be quasi-characters with exact
multiplicative conductors.  Suppose
\[
 m_F(\theta)=q=2d+\varepsilon>1,
 \qquad \varepsilon\in\{0,1\},
 \qquad m_F(\nu)\leq d.
\]
Put $r=d+\varepsilon$.  Then $q\leq2r$ and $r+1\leq q$, so $r$ is the
minimal permitted Lamprecht depth.  Moreover $m_F(\nu)<q$, hence
$\nu\theta$ has exact conductor $q$, and $\nu$ is trivial on $U_F^r$.

Choose an admissible denominator $\Gamma\in F^\times$, so
\[
 \operatorname{ord}_F(\Gamma)=q+n.
\]
The stationary numerators at depth $r$ are the quotient classes
\[
 \operatorname{St}_{\Gamma,r}(\theta),
 \operatorname{St}_{\Gamma,r}(\nu\theta)
 \in \mathcal O_F/\mathfrak p_F^{q-r}
     =\mathcal O_F/\mathfrak p_F^d.
\]
They are equal:
\[
 \operatorname{St}_{\Gamma,r}(\nu\theta)
 =\operatorname{St}_{\Gamma,r}(\theta).
\]
This is an equality of quotient classes, not a choice of a canonical
representative.  Indeed, every representative $\beta$ of the class of
$\theta$ satisfies
\[
 \theta(1+x)=\psi(\beta x/\Gamma)
 \qquad(x\in\mathfrak p_F^r),
\]
and multiplication by $\nu(1+x)=1$ gives the same identity for
$\nu\theta$.  Exact stationary order proves
$\operatorname{ord}_F(\beta)=0$, so each such $\beta$ is a nonzero local
unit and $\Gamma/\beta$ is defined.

For every representative $\beta$ of this stationary class, the finite
Gauss sums satisfy the exact identity
\[
 G_\Gamma(\nu\theta,\psi)
 =\nu(\beta)^{-1}G_\Gamma(\theta,\psi).
\]
Consequently
\[
 \boxed{
 \Delta_F^{\mathrm{fin}}(\nu\theta,\psi;\Gamma)
 =\nu(\Gamma/\beta)\,
  \Delta_F^{\mathrm{fin}}(\theta,\psi;\Gamma).}
\]
The quotient factor is representative-independent: if $\beta'$ is another
representative, then
\[
 \beta'/\beta\in U_F^d,
 \qquad
 \nu(\beta'/\beta)=1,
 \qquad
 \nu(\Gamma/\beta')=\nu(\Gamma/\beta).
\]
\LeanFile{LanglandsFirstMainLemma/Lamprecht/StableTwist.lean}
\lean{LanglandsFirstMainLemma.stableTwistData}
\lean{LanglandsFirstMainLemma.stableTwistData_conductor}
\lean{LanglandsFirstMainLemma.stableTwist_stationaryDepth}
\lean{LanglandsFirstMainLemma.stableTwist_trivialOnStationaryLayer}
\lean{LanglandsFirstMainLemma.stationaryNumeratorClass_stableTwist}
\lean{LanglandsFirstMainLemma.stableStationaryRepresentativeUnit}
\lean{LanglandsFirstMainLemma.stableTwist_characterFactor_representative_independent}
\lean{LanglandsFirstMainLemma.stableTwistAdmissibleGamma}
\lean{LanglandsFirstMainLemma.finiteGaussSum_stableTwist}
\lean{LanglandsFirstMainLemma.stableTwist}
\leanok
\SourceText{Lemma~\ref{lem:stable-twist-new}.}
\uses{mod:lamprecht-formula,mod:lamprecht-stationaryclasscalculus,mod:delta-elementary}
\FormalizationContract The parity parameter in Lean is a natural number
$\varepsilon\leq1$, hence exactly zero or one, and the stationary depth is
literally $d+\varepsilon$.  The twist datum stores the unchanged exact
conductor $q$ and the same admissible denominator.  The stationary numerator
remains a class in the certified lattice quotient; every theorem about its
value quantifies over a supplied representative and its proof of membership
in that class.  The finite-sum proof averages over $U_F^r/U_F^q$ and uses the
perfect Lamprecht pairing to annihilate every orbit except
$u\equiv\beta\pmod{\mathfrak p_F^d}$.  It therefore treats the even branch
$r=d$ and odd branch $r=d+1$ with their exact common denominator depth $d$,
without importing a canonical representative or discarding an odd residual
phase.  The principal theorem has the factor $\nu(\Gamma/\beta)$, with the
inverse on $\beta$ and no additional sign.
\end{theorem}
\begin{proof}
Since $m_F(\nu)\leq d<q=m_F(\theta)$, the unequal-conductor product lemma
gives $m_F(\nu\theta)=q$.  The inequalities
$q\leq2(d+\varepsilon)$ and $d+\varepsilon+1\leq q$ follow from
$q=2d+\varepsilon>1$ and $\varepsilon\leq1$.  Thus both characters use the
same stationary quotient at depth $r=d+\varepsilon$.  The character $\nu$
is trivial on $U_F^r$, so stationary-duality uniqueness gives equality of
their stationary classes.

Fix an arbitrary representative $\beta$.  To prove the Gauss-sum identity
without selecting representatives of any finite quotient, write
$G=U_F^0/U_F^q$ and $H=U_F^r/U_F^q$.  For $u\in G$, let
\[
 f(u)=\psi(u/\Gamma)\theta(u)^{-1}.
\]
For $a\in H$, stationary linearization gives
$\theta(a)=\psi(\beta(a-1)/\Gamma)$, and therefore
\[
 f(au)
 =\psi\bigl((u-\beta)(a-1)/\Gamma\bigr)f(u).
\]
The character in $a$ is the Lamprecht pairing of the coefficient class
$[u-\beta]\in\mathcal O_F/\mathfrak p_F^d$ with
$[a-1]\in\mathfrak p_F^r/\mathfrak p_F^q$.  Perfectness of the pairing and
finite additive-character orthogonality show that its sum over $H$ is zero
unless $[u-\beta]=0$.  In that exceptional case
$u/\beta\in U_F^d$, so $\nu(u)=\nu(\beta)$.

Now average
\[
 \sum_{u\in G}
 \bigl(\nu(u)^{-1}-\nu(\beta)^{-1}\bigr)f(u)
\]
over multiplication by $H$.  The factor in parentheses is constant on each
$H$-orbit.  The preceding orthogonality kills every nonexceptional orbit,
and it is zero on every exceptional orbit.  Since $H$ is nonempty, the
original sum is zero.  This proves
$G_\Gamma(\nu\theta,\psi)=\nu(\beta)^{-1}G_\Gamma(\theta,\psi)$.

The finite Gauss sum for $\theta$ is nonzero, and $\nu(\beta)$ has complex
norm one because $\beta$ is a local unit.  Hence taking total phases preserves
the factor $\nu(\beta)^{-1}$.  Multiplying by the denominator value in the
definition of $\Delta_F^{\mathrm{fin}}$ gives
\[
 (\nu\theta)(\Gamma)\nu(\beta)^{-1}
 =\nu(\Gamma/\beta)\theta(\Gamma),
\]
which is the displayed stable-twist formula with the required direction.

Finally, two representatives differ modulo $\mathfrak p_F^d$ and both have
order zero.  Thus their ratio lies in $U_F^d$, where $\nu$ is trivial, proving
the representative-independence assertion.
\end{proof}
